\section{Conclusion}
\label{sec:conclusion}

This paper introduced CodeGraph, a hybrid retrieval system for AI coding agents that bridges the gap between natural-language task descriptions and repository-level source code. By extracting lexical seeds and propagating relevance over a static dependency graph using Personalized PageRank, CodeGraph achieves 74.0\% Recall@10 on the SWE-bench Lite benchmark.

The central finding of this work is that structural retrieval performance is primarily bottlenecked by Reachability, not Ranking. When lexical signals identify even a noisy starting point, the graph topology reliably arbitrates relevance and pushes correct files to the top. The Seed-First engineering strategy (incorporating uniform PPR restart, compound tokenization, and IDF edge weighting) outperformed the lexical baseline by 24.3 percentage points.

Future work should target the remaining Connectivity and Semantic Gaps. Dense embedding models could replace sparse matching to better map abstract language to starting nodes, while lightweight dynamic tracing could expose the implicit dependencies that static analysis misses. Ultimately, CodeGraph demonstrates that repository retrieval improves when text matching is constrained to finding starting points, while deterministic graph propagation identifies the code that truly matters.
